%%%%%%%%%%%%%%%%%%%%%%%%%%%%%%%%%%%%%%%%%%%%%%%%%%%%%%%%%%%%%%%%%%%%%%
% Lớp 11 | Quan hệ vuông góc trong không gian | Phép chiếu vuông góc
% Dạng 1H8?3-1: Lý thuyết về phép chiếu vuông góc
% Số câu: 1  (sinh tự động bởi tach_cau.py; từ nay thầy sửa trực tiếp file này)
%%%%%%%%%%%%%%%%%%%%%%%%%%%%%%%%%%%%%%%%%%%%%%%%%%%%%%%%%%%%%%%%%%%%%%

%%%%%%%%%%%%%%%% TRẮC NGHIỆM 4 PHƯƠNG ÁN %%%%%%%%%%%%%%%%

% nguồn: Ok_TracNghiemLop11_Lan1
\begin{ex}%[1H8H3-1]%[ma:f2351e51]
	Cho hình chóp  $S.ABCD$ có đáy là hình vuông tâm $O$ và $SA=SB=SC=SD$. Hình chiếu của điểm $S$ trên mặt phẳng $(ABCD)$ là điểm nào dưới đây?
	 \choice
		{\True  $O$}
		 {$A$}
		{$D$} 
		 {$C$}
	 \loigiai{
		\begin{center}
			\begin{tikzpicture}[declare function={r=3; gd=-147;},font=\scriptsize]
					\path  (0,0) coordinate (A)
						 (0:r) coordinate (D)
						(gd:0.6*r)  coordinate (B)
						($(B)+(D)-(A)$)  coordinate (C)
						($(B)!0.5!(D)$) coordinate  (O)
						($(O)+(90:1.2*r)$) coordinate (S); 
					\draw (S)--(D)--(C)--(B)--cycle--(C); 
					\draw[dashed]  (S)--(A)--(B)--(D)--(A)--(C);
					 \draw[dashed] (S)--(O);
					\foreach \t/\g  in {A/140,B/180,C/-60,D/0,S/90,O/-90}{
						\fill (\t) circle (1pt)  node[shift={(\g:7pt)},font=\scriptsize]{$ \t $}; }
				\end{tikzpicture}
		\end{center}
		Vì $SA=SB=SC=SD$ nên hình chiếu của điểm $S$ trên mặt phẳng $(ABCD)$ là $O$ tâm  đường tròn ngoại tiếp đáy.
	} 
\end{ex}
